\subsection*{5.2 A reduction from Graph Coloring to Partial-MCCP}
Let us be given a graph $G=([n],E)$ with $E\subseteq\binom{[n]}2$, and consider the partial two-player Boolean function $f_G:[n]\times[n]\to\{0,1,*\}$, given by
\[f_G(i,j)=\begin{cases}1&\text{if }i=j,\\0&\text{if }\{i,j\}\in E,\\ *&\text{if }\{i,j\}\notin E.\end{cases}\]
By $f_G(i,j)=*$ we mean that $f_G(i,j)$ is undefined, which is to say Alice and Bob may output any value when given $(i,j)$ as input.
\setcounter{nttheorem}{31}
\begin{nttheorem}[Author Theorem32]
We have $L(f_G)\leq2\chi(G)$ and $C_1(f_G)=\chi(G)$.
\end{nttheorem}
\begin{proof}
To show that $L(f_G)\leq2\chi(G)$, consider the simple protocol where Alice and Bob have agreed on a coloring of $G$, so Alice begins by sending the color of her vertex, and Bob replies whether the color of his vertex is the same. If the two colors match, they output 1, and otherwise they output 0.
This is a protocol for $f$, since a valid coloring of $G$ will only color $i$ and $j$ by the same color if $i=j$, or $\{i,j\}\notin E$, in which case 1 is a valid output. If $i$ and $j$ are colored differently, then $i\ne j$, and so 0 is a valid output. This protocol has $2\chi(G)$ leaves.
To show that $C_1(f_G)\leq\chi(G)$, we cover the diagonal by 1-monochromatic rectangles; each such rectangle corresponds to a color $c$, and is the smallest rectangle containing all diagonal entries $(i,i)$ for vertices $i$ of $G$ that are colored $c$.
Now suppose we have any positive cover $\Gamma$ for $f$. The diagonal entries must be covered by 1-monochromatic rectangles of $\Gamma$, so consider the coloring of $G$ which colors vertex $i\in[n]$ by the 1-monochromatic rectangle which covers $(i,i)$. By our choice of $f$, any two diagonal entries belonging to the same 1-monochromatic rectangle are not connected by an edge of $G$. And so this gives us a valid coloring of $G$, which implies that $\chi(G)\leq|\Gamma|$.
\end{proof}
\subsection*{5.3 Consequences of the reduction}
We may now take Theorems32 and20, together with Lemmas28 and31, to obtain:
\begin{ntcorollary}[Author Corollary33]
For any constant $\epsilon>0$, it is NP-hard to approximate $L(f),P(f),C(f),C_1(f)$ and $S(f)$, for a given partial function $f:[n]\times[n]\to\{0,1,*\}$, with an approximation ratio better than $n^{1-\epsilon}$.
\end{ntcorollary}
Observe that this hardness of approximation result with a ratio of $n^{1-\epsilon}$ is very close to optimal, since all aforementioned measures for a partial two-player function $f:[n]\times[n]\to\{0,1,*\}$ are upper-bounded by $2n$. Improving upon this might still be possible, and we leave it as an open problem.
Now notice that the protocol given in the proof of Theorem20 is balanced. For this reason, its depth is logarithmic in the number of leaves. Thus, since computing $L(f)$ is hard up to an approximation ratio of $n^{1-\epsilon}$, it follows that computing the communication complexity $D(f)$ for a given partial function $f:[n]\times[n]\to\{0,1,*\}$ is hard, even if we allow for an additive error term of $(1-\epsilon)\log n$. Since $D(f)\leq\log n$, this implies that we cannot approximate $D(f)$ with an approximation ratio better than $\approx1/\epsilon$, which we can take as an arbitrarily large constant. The same reasoning applies to non-deterministic communication complexity.
\begin{ntcorollary}[Author Corollary34]
For any constant $\epsilon>0$, it is NP-hard to approximate $D(f)$ or $N(f)$, for a given partial function $f:[n]\times[n]\to\{0,1,*\}$, with an error term smaller than $(1-\epsilon)\log n$. For any constant $c>1$, it is NP-hard to approximate $D(f)$ or $N(f)$, for a given partial function $f:[n]\times[n]\to\{0,1,*\}$, with an approximation ratio better than $c$.
\end{ntcorollary}
Let us now prove that, assuming $\operatorname P\ne\operatorname{NP}$, there is no polynomial-time computable approximate characterization of the communication complexity of a partial Boolean function. This is captured by a more general result (Theorem35, stated at the start of this item), which is an immediate consequence of the reduction presented in Section5.2 and Theorem21.
\begin{proof}[Author proof of Theorem35]
Suppose that functions $\lambda,\eta$, and $r$ exist as described in the statement of the theorem. Let $L$ be any language in $\operatorname{NP}$, and assume that $k$ is a sufficiently large constant. For any $x\in\{0,1\}^n$, let $G_x$ be the corresponding graph given by Theorem21. Then let $f_x=f_{G_x}$ be as defined in Section5.2 above, and let $M_x$ be the communication matrix of $f_x$. If $x\in L$, we then have by Theorem32 that $k\geq\chi(G_x)\geq\frac12L(f_x)$, and because the protocol for $f_x$ that witnesses this fact is balanced, we get $1+\log k\geq D(f_x)$. Hence $\eta(1+\log k)\geq\eta(D(f_x))\geq r(M_x)$. If, on the other hand,$y\notin L$, then we have $2^{0.99k}<\binom k{\lfloor k/2\rfloor}<\chi(G_y)=C_1(f_y)\leq L(f_y)$, and thus $0.99k<D(f_y)$. Then $\lambda(0.99k)\leq\lambda(D(f_y))\leq r(M_y)$. But since $\lambda(0.99k)>\eta(1+\log k)$ using our assumptions on these functions and on $k$, it follows that $r(M_y)>r(M_x)$. As a consequence, the polynomial time function $r$ can be used to distinguish positive and negative instances of $L$. Since $L$ is an arbitrary language in $\operatorname{NP}$, we get that $\operatorname P=\operatorname{NP}$. This completes the proof.
\end{proof}
